\chapter{Discussion}

The empirical analysis presented in the previous chapter yielded several significant insights into DAO activity measurement, while also challenging some conventional assumptions about organizational assessment in decentralized contexts. This chapter examines these findings' theoretical and practical implications, while also identifying promising directions for future research.

\section{Theoretical Implications}

Our investigation into DAO activity measurement contributes to management theory in several ways. Most notably, our findings extend \textcite{pfeffer2015external}'s resource dependence theory (RDT) in the context of DAOs. While RDT traditionally emphasizes financial resources as critical for organizational survival, our results suggest that for DAOs, development infrastructure maintenance represents an equally crucial resource. This insight is particularly relevant for technology-driven organizations and non profit organizations, where continuous technical development may prove more vital than financial metrics for organizational sustainability.

The study also advances our understanding of digital governance effectiveness measurement. Previous research has primarily focused on governance structure design without providing empirical methods for measuring ongoing effectiveness (\cite{chohan2017decentralized}; \cite{wang2019decentralized}; \cite{rikken2019governance}). Our LSTM framework offers a novel quantitative approach to this challenge, providing a systematic method for evaluating governance effectiveness through multiple complementary metrics. This addresses a significant gap in the literature regarding empirical assessment methods for decentralized organizations (\cite{tan2023open}).

\section{Methodological Innovation}

The LSTM model developed in this study represents a significant methodological contribution to organizational research. While our findings indicate that development activity can serve as a reliable proxy for DAO activity in many cases, the LSTM framework's sophistication offers several distinct advantages. Unlike point-in-time metrics, our model captures the dynamic evolution of organizational behavior over time, enabling more nuanced activity classification through temporal pattern recognition (\cite{wang2003temporal}). Additionally, the model synthesizes financial, social, and technical indicators, providing a foundation for incorporating additional metrics as the field evolves through multi-dimensional integration. Finally, although the current implementation achieves 75\% accuracy with just three input dimensions, the framework demonstrates considerable potential for enhancement through additional data sources and refined architecture, highlighting its predictive potential.












\section{Practical Implications}

Our research findings yield substantial practical implications for key stakeholders in the DAO ecosystem. Firstly, for DAO operators, our results indicate the critical importance of maintaining consistent development activity throughout market cycles. This finding challenges the common practice of scaling technical operations based on market conditions (\cite{teece1997dynamic}), suggesting instead that sustained development effort serves as a fundamental indicator of organizational vitality. 

For institutional investors, our findings offer evidence-based metrics for evaluating DAO sustainability. The research indicates that while token price movements and social media engagement provide immediate market signals (\cite{chen2009approximability}), development activity trends offer more reliable indicators of long-term organizational health. This insight enables more sophisticated due diligence processes that extend beyond traditional market metrics.

The demonstrated effectiveness of the LSTM model in predicting DAO activity states provides valuable insights for platform operators and information aggregators. Current platforms predominantly emphasize financial and social metrics (\cite{deepDaoWebsite}), potentially misaligning user attention with indicators of organizational sustainability. Our findings suggest that platforms should consider elevating the visibility of development activity metrics while maintaining a balanced presentation of all organizational indicators.

The research also carries implications for regulatory frameworks surrounding decentralized organizations. The ability to quantifiably measure DAO activity through our model provides a foundation for developing standardized reporting requirements. Importantly, the effectiveness of simple activity metrics suggests the potential for regulatory approaches that achieve oversight objectives without imposing undue operational burdens on decentralized organizations.

Of particular significance is the finding that continuous monitoring of activity patterns provides more valuable insights than point-in-time assessments. This observation challenges traditional organizational evaluation approaches and suggests the need for dynamic assessment frameworks that can capture the evolving nature of decentralized operations. The temporal nature of our LSTM model provides a promising template for such frameworks, offering a bridge between the need for ongoing oversight and the reality of decentralized governance structures.










\section{Future Research Directions}

Our findings open several promising avenues for future research in decentralized organizational dynamics. A particularly compelling direction involves comparative analysis across different types of DAOs. While our current study establishes foundational metrics for activity assessment, systematic investigation of how these patterns manifest across DeFi, GameFi, and Social DAOs could yield valuable insights into sector-specific indicators of organizational health.
The relationship between market shocks and DAO activity patterns presents another rich area for investigation (\cite{faqir2021effect}). The cryptocurrency ecosystem's susceptibility to external disruptions provides a unique laboratory for studying organizational resilience. Future research could examine how different types of shocks - whether technical failures, market crashes, or regulatory changes - impact DAO activity patterns. Such analysis could help identify characteristics that enable some organizations to maintain operational continuity while others falter.
Of particular scholarly interest is the potential for investigating recovery patterns in DAOs that successfully navigate periods of reduced activity. Understanding the mechanisms through which organizations restore operational momentum could provide valuable insights for both theoretical frameworks and practical governance. This research direction would benefit from longitudinal studies tracking DAOs from inception through various lifecycle stages, with particular attention to critical transition points and recovery episodes.

\section{Limitations}

Several limitations of our study warrant consideration and provide additional context for future research. Firstly, our reliance on Santiment as a single data source, while enabling consistent analysis, may introduce systematic biases in our findings. While our results demonstrate robust predictive capability, validation across multiple data sources would strengthen their generalizability.
The temporal scope of our analysis, while sufficient for identifying activity patterns, may not capture longer-term organizational dynamics. Extended longitudinal studies could reveal additional patterns or challenge current assumptions about DAO lifecycle progression. Furthermore, our focus on three primary metrics, while empirically justified, may not capture all relevant dimensions of organizational activity.

% Despite these limitations, our findings provide a robust foundation for understanding DAO activity patterns and establish a methodological framework for future investigations. The demonstrated effectiveness of our LSTM model, combined with the validation of development activity as a key indicator, opens new possibilities for both theoretical advancement and practical application in the field of decentralized organizational research.




This chapter has examined how our investigation into DAO activity measurement advances both management theory and practice. Our analysis demonstrates that while development activity serves as an effective predictor of organizational vitality, the LSTM framework establishes groundwork for more sophisticated activity prediction approaches. The research contributes methodologically through its novel application of temporal pattern recognition, providing researchers and practitioners with a flexible toolkit for organizational assessment. While our findings challenge traditional assumptions about organizational measurement, particularly regarding financial and social metrics, they reinforce the fundamental importance of maintaining core operational capabilities. These insights establish a foundation for future investigations into decentralized organizational dynamics, particularly in understanding how different types of DAOs maintain operational continuity through varying market conditions.






